% Source: https://artofproblemsolving.com/community/c3h1987943_algebra_help
% Stable identity: aops-topic-1987943-q01
Let $A=a+\sqrt{a^2+1}>0$. Then $\dfrac{1}{A}=\sqrt{a^2+1}-a$ (multiply: $(a^2+1)-a^2=1$), so
\[
\sqrt{a^2+1}=\dfrac{1}{2}\left(A+\dfrac{1}{A}\right),\qquad a=\dfrac{1}{2}\left(A-\dfrac{1}{A}\right).
\]
Similarly, with $B=b+\sqrt{b^2+1}>0$, $\sqrt{b^2+1}=\dfrac{1}{2}\left(B+\dfrac{1}{B}\right)$ and $b=\dfrac{1}{2}\left(B-\dfrac{1}{B}\right)$. Since
\[
\dfrac{1}{A}+\dfrac{1}{B}=\dfrac{A+B}{AB},
\]
we get
\[
\sqrt{a^2+1}+\sqrt{b^2+1}=\dfrac{1}{2}(A+B)\left(1+\dfrac{1}{AB}\right),\qquad a+b=\dfrac{1}{2}(A+B)\left(1-\dfrac{1}{AB}\right).
\]
Since $AB=2019\ne1$ and $A+B>0$, the denominator $a+b$ is nonzero, and
\[
\dfrac{\sqrt{a^2+1}+\sqrt{b^2+1}}{a+b}=\dfrac{AB+1}{AB-1}=\dfrac{2020}{2018}=\boxed{\dfrac{1010}{1009}} .
\]
